\documentclass[a4paper,11pt]{report}

\usepackage[utf8]{inputenc}
\usepackage[T1]{fontenc}
\usepackage{array}
\usepackage{longtable}
\usepackage{hyperref}

\newcommand{\mr}{\textsc{MR}}
\newcommand{\mrmac}{\textsc{MRMAC}}
\emergencystretch=3em

\title{\mr\\Multi-Edit Revisited\\Userhandbuch v1.0}
\author{Projekt: \url{https://github.com/ebeneezer/MR}}
\date{22. Juni 2026}

\begin{document}
\maketitle

\begin{quote}
``Software and cathedrals are much the same -- first we build them, then we pray.''

\hfill Sam Redwine
\end{quote}

\tableofcontents
\listoftables

\chapter{Einfuehrung}

\section{Was ist MR?}

\mr{} ist eine Neuimplementierung des klassischen Programmier-Editors
Multi-Edit fuer moderne Linux-Terminals. Der Editor verbindet klassische
TUI-Bedienung mit heutigen Erwartungen an grosse Dateien, Syntaxanalyse,
Language-Server-Integration, Projektprofile und reproduzierbare
Einstellungen.

Der Schwerpunkt liegt nicht auf einer IDE-Oberflaeche, sondern auf einem
schnellen, tastaturzentrierten Arbeitsfluss. Fenster, Makros, Suche,
Compiler-Ausgaben und SideKicks bleiben im Terminal und sollen auch bei sehr
grossen Dateien reaktionsfaehig bleiben.

\section{Kernphilosophie}

\mr{} orientiert sich an drei praktischen Regeln:

\begin{itemize}
\item Wiederholbare Arbeit gehoert in Makros.
\item Editorfunktionen muessen auch bei grossen Dateien berechenbar bleiben.
\item Komfortfunktionen duerfen die direkte Bedienbarkeit nicht verdecken.
\end{itemize}

\section{Hauptfeatures auf einen Blick}

\begin{itemize}
\item eigene Makrosprache \mrmac{} mit Compiler und Coprozessor,
\item Piecetable-basierte Bearbeitung grosser Dateien,
\item Syntax-Highlighting, Folding, Outline und Smart-Indenting,
\item Language-Server-Protocol fuer externe Sprachdienste,
\item Compilerprofile mit Problems-Navigation,
\item virtuelle Desktops, Workspaces und Farbthemen.
\end{itemize}

\chapter{Installation}

\section{Systemvoraussetzungen}

\begin{table}[h]
\centering
\begin{tabular}{>{\bfseries}p{0.28\textwidth}p{0.58\textwidth}}
Betriebssystem & Linux, 64 Bit empfohlen\\
Terminal & UTF-8-faehig mit ncursesw-Unterstuetzung\\
Compiler & C++20-faehiger Compiler, im Projekt bevorzugt \texttt{clang++}\\
Bibliotheken & ncursesw, PCRE2, TVision im Projektkontext\\
\end{tabular}
\caption{Minimale Voraussetzungen}
\end{table}

\section{Installation aus dem Quellcode}

\begin{verbatim}
git clone https://github.com/ebeneezer/MR
cd MR
make clean all CXX=clang++
\end{verbatim}

\section{Erster Start}

Nach dem Build kann \mr{} direkt aus dem Arbeitsbaum gestartet werden. Die
Einstellungen werden ueber die vorgesehenen Setup- und Runtime-Pfade geladen;
Settings-Dateien sollen nicht durch externe Werkzeuge am Editor vorbei
synthetisiert werden.

\chapter{Die MRMAC-Makrosprache}

\section{Einfuehrung in MRMAC}

\mrmac{} ist die Makrosprache des Editors. Sie dient nicht nur der
Tastenbelegung, sondern auch der Automatisierung wiederkehrender
Editorablaeufe. Makros werden kompiliert und danach durch die VM beziehungsweise
den Coprozessor ausgefuehrt.

\subsection{Der Makro-Compiler}

Der Makro-Compiler uebersetzt Makroquellen in eine Form, die zur Laufzeit
ausgefuehrt werden kann. Fehler sollen dabei frueh sichtbar werden, damit
fehlerhafte Makros nicht erst in einer interaktiven Sitzung auffallen.

\subsection{Der Makro-Coprozessor}

Lang laufende oder deferred UI-nahe Operationen laufen nicht als lose
Nebenlaeufigkeit durch die Anwendung. Der Coprozessor kapselt Ausfuehrung,
Status und Rueckmeldung so, dass die VM-Grenzen nachvollziehbar bleiben.

\subsection{Beispiel: Einfaches MRMAC-Makro}

\begin{verbatim}
PROC HELLO_MR
  UI_MESSAGEBOX "MR" "Hello from MRMAC"
ENDPROC
\end{verbatim}

\section{Makro-Manager}

Der Makro-Manager macht Makros auffindbar und ausfuehrbar. Fuer produktive
Konfigurationen ist wichtig, dass die Runtime-Pfade und die persistierten
Settings konsistent bleiben.

\chapter{Editieren wie ein Profi}

\section{Umgang mit grossen Dateien}

\mr{} ist auf sehr grosse Dateien ausgelegt. Die interne Textdarstellung
vermeidet unnoetige Vollkopien und erlaubt Navigation, Suche und Anzeige auch
dann, wenn die Datei deutlich groesser als ein typischer Quelltext ist.

\section{Blocktypen}

Der Editor unterscheidet klassische Blockoperationen, Spaltenbloecke und
zwischen Fenstern nutzbare Kopierpfade. Blockoperationen sollen vorhersehbar
bleiben, auch wenn freie Cursorpositionen oder unterschiedliche Tabulatorbreiten
im Spiel sind.

\section{Code-Faltung und Smart-Indenting}

Folding und Smart-Indenting werden sprachspezifisch berechnet. Fuer LaTeX sind
insbesondere Umgebungen, Sections und Rohtextbloecke wie \texttt{verbatim}
wichtig:

\begin{verbatim}
\section{Dies ist hier nur Text}
\begin{itemize}
  \item Auch dies bleibt Rohtext.
\end{itemize}
\end{verbatim}

Der Editor darf diese innere Section nicht als echte Dokumentstruktur werten.

\section{Die Minimap}

Die Minimap ist eine kompakte Navigationshilfe. Sie ersetzt nicht die Outline,
sondern ergaenzt die schnelle Orientierung in langen Dateien.

\chapter{Fenster-Management}

\section{Der Fenster-Manager}

Fenster sind ein zentraler Bestandteil des Arbeitsflusses. Editorfenster,
SideKicks, Buildausgaben und Dialoge muessen ihre Rollen klar behalten.

\section{Virtuelle Desktops}

Virtuelle Desktops gruppieren Arbeitskontexte. Sie sollen schnelle
Kontextwechsel erlauben, ohne Fensterzustand oder Projektbezug zu verlieren.

\section{Inter-Window Copy \& Paste}

Kopieren zwischen Fenstern ist ein eigener Pfad und muss den Quell- und Zieltyp
korrekt beachten. Das gilt besonders fuer Spaltenbloecke und freie
Cursorpositionen.

\chapter{Workspaces und Projekte}

\section{Workspace-Management}

Workspaces speichern Arbeitskontexte. Sie sind settingsnah, aber nicht
automatisch Teil des kanonischen Kernsettings-Vertrags.

\section{Profile}

Profile bestimmen Sprache, Compiler und Editorverhalten fuer Dateitypen.
Die Zuordnung muss nachvollziehbar sein, damit ein geoeffnetes Dokument nicht
scheinbar zufaellig andere Editorregeln verwendet.

\section{Farbthemen}

Farbthemen betreffen die Darstellung, nicht die Semantik. Sie duerfen deshalb
keine verdeckte Abhaengigkeit fuer Parser, Build oder Runtime-Logik bilden.

\chapter{Suche und Ersetzung}

\section{Regulaere Ausdruecke}

\mr{} verwendet PCRE2 fuer regulaere Ausdruecke. Suchen sollen schnell bleiben
und Treffer so markieren, dass Navigation und Wiederholung stabil sind.

\section{Rekursive Multifile-Suche}

Die rekursive Suche ist fuer Projektarbeit gedacht. Trefferlisten muessen
bedienbar bleiben und duerfen den Editor nicht blockieren.

\chapter{Compiler-Integration}

\section{Compiler-Profile}

Compilerprofile beschreiben, wie ein Dokument oder Projekt gebaut wird. Fuer
LaTeX ist \texttt{latexmk} der Standardpfad:

\begin{verbatim}
latexmk -pdf documentation/manuals/mr-userref.tex
\end{verbatim}

\section{Fehler-Tracking}

Build-Ausgaben werden in Problems ueberfuehrt, wenn sie verwertbare
Quellpositionen oder eindeutige Fehlermeldungen enthalten. LaTeX-Fehlerzeilen
mit fuehrendem Ausrufezeichen und Warnungen mit \texttt{on input line N} sind
fuer diesen Pfad wichtig.

\chapter{Tastatur und Keymapping}

\section{Key-Mapping-Manager}

Der Key-Mapping-Manager ordnet Tastenfolgen Editorfunktionen oder Makros zu.
Resolver-Semantik und Persistenz sind geschuetzte Architekturteile.

\section{Emulationsmodi}

Emulationsmodi erleichtern Umstieg und Muskelgedaechtnis. Sie duerfen aber
nicht dazu fuehren, dass interne Hotkeys hart codiert an UI-Texte gekoppelt
werden.

\chapter{Erweiterte Funktionen}

\section{Drucken und Export}

Druck- und Exportpfade sollen reproduzierbar sein. Fuer dieses Handbuch ist der
LaTeX-PDF-Pfad zugleich ein Akzeptanzfall fuer Editor, Build und Diagnostics.

\section{Pipes und Shell-Kommandos}

Shell-Kommandos laufen ueber die konfigurierte Shell. Direkte Rueckfaelle auf
\texttt{system()} oder \texttt{popen()} waeren fuer diesen Pfad falsch.

\section{Automatische Sprachdetektion}

Die Sprachdetektion verbindet Dateiendung, Profil und Syntaxregeln. Fuer
\texttt{.tex}-Dateien muessen LaTeX-Regeln gelten, nicht zufaellig passende
C- oder C++-Heuristiken.

\chapter{Tastaturreferenz}

\begin{longtable}{>{\ttfamily}p{0.24\textwidth}p{0.62\textwidth}}
Alt-Enter & Snippet-SideKick-Auswahl uebernehmen\\
Tab & naechsten Snippet-Platzhalter ansteuern\\
Shift-Tab & vorherigen Snippet-Platzhalter ansteuern\\
F4 & Suchtreffer beziehungsweise markierte Position fortsetzen\\
\end{longtable}

\chapter{Fehlerbehebung}

\section{Haeufige Probleme}

Wenn ein LaTeX-Build fehlschlaegt, ist zuerst zwischen Paketfehler,
Syntaxfehler, Warnung und Mehrfachlauf-Hinweis zu unterscheiden. Nicht jede
LaTeX-Zeile ist ein navigierbares Problem, aber spezifische Fehlertexte muessen
sichtbar bleiben.

\chapter{Technische Architektur}

\section{Verwendete Technologien}

\mr{} nutzt C++20 als Compilerstandard, bleibt stilistisch aber bewusst nahe an
klassischem C++18. Externe Bibliotheken werden dort eingesetzt, wo sie einen
klaren technischen Zweck erfuellen.

\section{Der Coprozessor im Detail}

Der Coprozessor trennt lang laufende Arbeit, UI-Rueckmeldung und VM-Grenzen.
Diese Trennung ist fuer Buildausgaben, deferred UI und Makroausfuehrung
entscheidend.

\appendix

\chapter{MRMAC-Sprachreferenz}

\section{Grundlegende Syntax}

\mrmac{}-Programme bestehen aus Prozeduren, Anweisungen und Aufrufen von
Editorfunktionen. Die genaue Semantik ist Teil der Makro- und VM-Vertraege.

\section{Editor-Kommandos}

Editor-Kommandos duerfen nicht als lose globale Registry neben den vorhandenen
Dispatch- und Settingspfaden wachsen.

\chapter{Lizenz und Danksagungen}

\section{Zitate}

Dieses Handbuch enthaelt kurze historische Referenzen auf Multi-Edit und die
Motivation hinter \mr{}. Der Fokus der aktuellen Quelle liegt auf einem
editierbaren LaTeX-Akzeptanzdokument.

\section{Projekt}

Projektseite: \url{https://github.com/ebeneezer/MR}

\end{document}
